\textbf{Área de superfície:}
\[
A(R)=4\pi R^2.
\]
Para raio unitário: $A(1)=4\pi$. Logo:
\[
\frac{A(R)}{A(1)}=\frac{4\pi R^2}{4\pi}=R^2.
\]

\textbf{Volume:}
\[
V(R)=\frac{4}{3}\pi R^3.
\]
Para raio unitário: $V(1)=\frac{4}{3}\pi$. Logo:
\[
\frac{V(R)}{V(1)}=\frac{\tfrac{4}{3}\pi R^3}{\tfrac{4}{3}\pi}=R^3.
\]

\textbf{Verificação numérica} com $R=2$:
\[
A(2)=4\pi(4)=16\pi\approx50{,}27, \quad A(1)=4\pi\approx12{,}57,
\]
\[
\frac{A(2)}{A(1)}=4=2^2.\quad \checkmark
\]
\[
V(2)=\frac{4}{3}\pi(8)=\frac{32\pi}{3}\approx33{,}51, \quad V(1)=\frac{4\pi}{3}\approx4{,}19,
\]
\[
\frac{V(2)}{V(1)}=8=2^3.\quad\checkmark
\]

\textbf{Conclusão:} em geral, para um fator de escala $\lambda$ qualquer:
\[
A(\lambda R)=\lambda^2 A(R),\qquad V(\lambda R)=\lambda^3 V(R).
\]
A área escala com a \emph{segunda potência} do comprimento e o volume com a \emph{terceira potência}, consequência direta do caráter bidimensional e tridimensional dessas grandezas geométricas.

